\chapter{Descripción del proyecto}\label{cap:descripcion_del_proyecto}
\section{Introducción}
La Diabetes Mellitus es un trastorno metabólico crónico caracterizado por niveles elevados de glucosa en sangre debido a una deficiencia en la producción o acción de la insulina\cite{fede}. Aunque se describe desde el Antiguo Egipto en el Papiro de Ebers (1500 a. C.) como una enfermedad de ``orina dulce``, no fue hasta 1921 cuando Banting y Best transformaron su pronóstico \negritas{antes mortal} con el descubrimiento de la insulina. Desde entonces, la evolución ha pasado de la supervivencia básica a una gestión tecnológica avanzada basada en la monitorización digital y la precisión farmacológica.

Esta enfermedad ha experimentado un crecimiento alarmante en las últimas décadas. Su gestión diaria requiere que el paciente tome decisiones constantes y complejas sobre la administración de insulina, la ingesta de hidratos de carbono y la actividad física. A pesar de los avances en la tecnología médica, como los sistemas de monitorización continua de glucosa (MCG), el manejo de la patología sigue dependiendo en gran medida del conocimiento y la capacidad de interpretación del usuario.

En este trabajo de fin de grado se realizará el diseño y desarrollo de una aplicación móvil integral para la gestión avanzada de la diabetes, centrada en la convergencia de la tecnología NFC y la Inteligencia Artificial. El proyecto comprende desde la implementación de la lectura directa de sensores de glucosa intersticial hasta la creación de un sistema de soporte inteligente integrado. A través de esta herramienta, se busca proporcionar al usuario un entorno capaz de centralizar sus pautas médicas, registrar su actividad física e interactuar con una interfaz que traduzca datos complejos en recomendaciones pedagógicas, facilitando así un aprendizaje continuo y una toma de decisiones más precisa en su día a día.

\section{Motivación}
La motivación de este proyecto es \negritas{profundamente personal} y nace de mi propia experiencia diaria conviviendo con la diabetes. Como paciente, identifico de primera mano las carencias de las soluciones actuales: mientras que aplicaciones oficiales como la del sistema FreeStyle Libre se limitan a una captura de datos básica y lineal, por otro lado, las alternativas existentes en el mercado suelen presentar barreras de acceso debido a su coste o carecen de un equilibrio entre simplicidad, usabilidad y funcionalidad, factores esenciales para pacientes con pocos conocimientos técnicos o recién diagnosticados.. Esta realidad genera una frustración constante, ya que nos encontramos con herramientas que registran cifras, pero que no ofrecen un soporte real ante la incertidumbre del día a día.

Como usuario y paciente directamente afectado, mi impulso es crear la herramienta que considero que falta en el mercado: una plataforma que combine simplicidad, usabilidad y funcionalidad máxima. Mi objetivo es que cualquier persona con diabetes, independientemente de sus conocimientos tecnológicos o médicos, pueda entender su metabolismo y gestionar su enfermedad sin la carga cognitiva que supone interpretar datos aislados.

Desde \negritas{la perspectiva de la ingeniería}, el reto que me he propuesto es transformar esta necesidad vital en una solución tecnológica integral "Full Stack". El proyecto supone para mí un desafío de integración multidisciplinar: desde la interacción directa con el hardware del móvil mediante el chip NFC, hasta el diseño de una arquitectura completa que incluya Frontend, Backend y gestión de bases de datos, culminando con la integración de Inteligencia Artificial como motor de ayuda pedagógica. En definitiva, mi motivación reside en poner la ingeniería al servicio del paciente, facilitando una gestión de la diabetes más humana, inteligente y, sobre todo, accesible.

\section{Objetivos}
El objetivo principal de este trabajo es el desarrollo de una solución tecnológica integral para la monitorización y gestión de la Diabetes Mellitus, que combine la adquisición de datos en tiempo real mediante hardware móvil y el soporte pedagógico basado en Inteligencia Artificial.

Para alcanzar este objetivo general, se proponen los siguientes objetivos específicos:

\begin{itemize}
\item \negritas{Implementar la lectura directa de sensores de glucosa mediante el uso del chip NFC} del dispositivo móvil, permitiendo la captura de datos del sensor FreeStyle Libre 2 Plus de forma instantánea.

\item \negritas{Diseñar y desarrollar una arquitectura de software completa} (Full Stack) que incluya una base de datos local y remota para la gestión del perfil diabético, pautas médicas, registro de ingestas y actividad física.

\item \negritas{Integrar un sistema de Inteligencia Artificial capaz de actuar como asistente educativo}, analizando variables críticas como la Insulina Activa (IOB) y el impacto del ejercicio para ofrecer recomendaciones orientativas al usuario.

\item \negritas{Optimizar la experiencia de usuario (UX) mediante una interfaz centrada en la simplicidad y la usabilidad}, eliminando las barreras técnicas presentes en las soluciones actuales para facilitar el acceso a pacientes sin conocimientos avanzados.

\item \negritas{Garantizar la seguridad y el rigor informativo mediante la configuración de un entorno donde la IA actúe exclusivamente como soporte pedagógico}, reforzando siempre la importancia de las pautas médicas establecidas.

\end{itemize}
